% authority: science_draft
% object-id: TEX-V22-P4-T02-B2-POST-LINE-LOCK-ROUTE-SELECTION-V1
% task-id: RT-20260810-002
% agentjob-id: AJ-RT-20260810-002-001
\documentclass[11pt]{article}

\usepackage[margin=0.86in]{geometry}
\usepackage{amsmath,amssymb,amsthm}
\usepackage{booktabs,longtable,array,enumitem}
\usepackage[hidelinks]{hyperref}

\newtheorem{definition}{Definition}
\newtheorem{lemma}{Lemma}
\newtheorem{proposition}{Proposition}
\newenvironment{nonconclusion}{\par\smallskip\noindent\textbf{Non-conclusion.}\ }{\par\smallskip}

\newcommand{\Seq}{\mathcal S_{\mathrm{eq}}}
\newcommand{\RR}{\mathbb R}
\newcommand{\GammaInt}{\Gamma_{\cap}}
\newcommand{\Pprod}{P_{\Pi}}

\title{V22 P4--T02 B2 Post-Line-Lock Route Selection v1\\
\large A common-hyperbolicity-envelope packet without shared-line replay}
\author{Theoretical Continuation Selector Packet RT-20260810-002}
\date{2026-08-10}

\begin{document}
\sloppy
\maketitle

\begin{abstract}
The preceding P4--T02 packet proves that fixed finite P7 records and a bare
smooth source patch do not select a projective tangent line naturally under
the full source-presentation group.  It also proves that equality of
independently varied sector lines has empty interior.  This selector therefore
excludes the shared-line, diagonal-only, and undeclared-line-stabilizer route.
It compares four remaining dispositions and selects one materially distinct
constructive packet: build or precisely obstruct a source-covariant product
principal polynomial and the intersection of independently supplied sector
hyperbolicity cones, then test whether a source-operational quotient justifies
that mathematical envelope as a generalized causal structure.  A finite
product lemma and an explicit three-sector split show why the construction is
mathematically available without equality of sector characteristic
hypersurfaces.  They do not establish universal propagation, a physical cone,
a Lorentzian conformal class, an effective metric, B2 activation, D7 adequacy,
or P4--T03 authority.
\end{abstract}

\section{Authority and fixed obstruction}

The fixed result is
\[
 \texttt{OBST-V22-P4T02-B2-NATURAL-LINE-LOCK-001},
\]
and the local freeze is
\[
 \texttt{NDCL-V22-P4T02-B2-SHARED-TAU-SELECTOR}.
\]
The freeze covers a single shared vector field inserted into every sector,
diagonal-only variations, or an equivalent restriction of source
presentations to the stabilizer of an undeclared line.  It does not freeze all
B2 descriptors or prove future source extensions impossible.

This task is a \emph{selector}.  It chooses one next packet but does not
execute that packet.  The source-extension budget, B1 termination, B2
inactivity, descriptor incompleteness, D7 non-evaluation, P4--T03 lock, and all
protected authority boundaries remain fixed.

\section{Decision rule and route comparison}

The selection obeys four rules.  First, a non-promotional constructive packet
is preferred when one exists.  Second, a route must be materially distinct
from equality of independently varied sector leading lines.  Third, it must
stay at the \texttt{effective\_metric\_g\_eff} milestone and inside P4--T02.
Fourth, it must distinguish mathematical common-time admissibility from
universal matter propagation and physical causal interpretation.

\begin{longtable}{>{\raggedright\arraybackslash}p{0.21\textwidth}>{\raggedright\arraybackslash}p{0.18\textwidth}>{\raggedright\arraybackslash}p{0.49\textwidth}}
\toprule
Route & Disposition & Exact basis\\
\midrule
Replay shared line locking & Excluded & It is the locally frozen strategy under a new name.\\
Add one common line primitive & Not selected & It preloads the failed equality property and is not materially distinct, even though proposal-only definition would be unprotected.\\
Construct a common hyperbolicity envelope & Selected & Product and cone intersection are natural operations on independently supplied sector principal data and do not require equal characteristic hypersurfaces.\\
Prove a broader B2 no-go & Not selected & A new constructive route is available, so a broader negative packet is premature.\\
Request ontology adoption or Gate review & Not selected & The construction can first be defined and refuted as draft/control without protected authority.\\
\bottomrule
\end{longtable}

The exact selected future packet is
\[
 \texttt{CAND-V22-B2-COMMON-HYPERBOLICITY-ENVELOPE-V1},
\]
assigned, after this task's checkpoint, to
\texttt{candidate-constructor@0.2.0} with embedded Refuter stress.

\section{Mathematical target}

Let \(U\) be a declared smooth source patch and let \(\Seq\) be a finite,
honestly enumerated sector domain.  For each \(s\in\Seq\), suppose a later
packet supplies a nonzero real homogeneous reduced principal polynomial
\[
 p_s(x,k),\qquad k\in T_x^*U,
\]
with declared degree, density weight, regularity, quotient provenance, and
source-presentation covariance.  These polynomials are not derived by the
present selector.  Under the current evidence they remain proposal-only
source-extension data.

\begin{definition}[Common hyperbolicity envelope]
Suppose every \(p_s\) is hyperbolic with respect to one covector \(h\), with
signs oriented so \(p_s(h)>0\).  Let \(\Gamma_s(h)\) be the hyperbolicity-cone
component of \(p_s\) containing \(h\).  Define
\[
 \Pprod(k)=\prod_{s\in\Seq}p_s(k),
 \qquad
 \GammaInt(h)=\bigcap_{s\in\Seq}\Gamma_s(h).
\]
The pair \((\Pprod,\GammaInt)\) is the proposed mathematical envelope.
\end{definition}

\begin{lemma}[Finite product and common component]
\label{lem:product}
For a finite family of real homogeneous polynomials hyperbolic with respect to
the same \(h\), \(\Pprod\) is hyperbolic with respect to \(h\).  Its
hyperbolicity-cone component containing \(h\) is \(\GammaInt(h)\).
\end{lemma}

The task-local identifier for this lemma is
\[
 \texttt{LEM-V22-P4T02-B2-COMMON-HYPERBOLICITY-PRODUCT-001}.
\]

\begin{proof}
For every covector \(q\), the roots of
\(t\mapsto \Pprod(q+t h)\) are the multiset union of the roots of the factors
\(t\mapsto p_s(q+t h)\).  All are real by the hypotheses, so the product is
hyperbolic.  The connected positive component containing \(h\) consists
exactly of covectors remaining in the corresponding positive component of
every factor; hence it is the intersection of the \(\Gamma_s(h)\).
\end{proof}

The construction is source-presentation covariant once the \(p_s\) are.  A
source diffeomorphism pulls back every sector polynomial, preserves products,
and transports intersections of their cone components.  No target metric,
coframe, tetrad, desired Lorentzian cone, or benchmark value is needed for
this algebraic operation.

\section{Explicit independent-sector witness}

At one source point choose a presentation basis \(e_0,e_1,e_2,e_3\) only for
calculation and set
\[
 v_R=e_0+\epsilon e_1,\qquad
 v_S=e_0,\qquad
 v_D=e_0-\epsilon e_1,
 \qquad 0<|\epsilon|<1.
\]
The three independently split projective lines are unequal.  Their linear
principal forms are
\[
 p_R(k)=k_0+\epsilon k_1,\quad
 p_S(k)=k_0,\quad
 p_D(k)=k_0-\epsilon k_1.
\]
With \(h=e^0\), every \(p_s(h)=1\), and
\[
 \Pprod(k)=k_0^3-\epsilon^2 k_0k_1^2.
\]
For any \(q\), the three roots of
\(t\mapsto\Pprod(q+t h)\) are
\[
 -q(v_R),\qquad -q(v_S),\qquad -q(v_D),
\]
and are real.  The common component is the nonempty open convex intersection
\[
 \GammaInt=\left\{\,k : k(v_R)>0,\;k(v_S)>0,\;k(v_D)>0\,\right\}.
\]
Strict positivity at \(h\) persists under sufficiently small independent
compact-open perturbations of the finite vector family.  Thus equality of
sector lines is not required for a robust common time-covector domain.

This witness is a selector-level feasibility proof only.  The individual
sector vectors are proposal-only data; they are not derived from the fixed P7
records.  The exact current protocol bridge remains restricted to the three
equipped controls and is sample-forgetting unless a finite preparation-set
extension is declared.

\section{The decisive remaining distinction}

The construction yields three different mathematical layers:
\begin{enumerate}[leftmargin=*]
  \item each sector has its own characteristic hypersurface
        \(\{p_s=0\}\);
  \item their product has the union of those hypersurfaces as its zero set;
  \item their oriented cone intersection gives a common admissible
        time-covector domain when nonempty.
\end{enumerate}
None of these facts alone says that every matter sector propagates on one
physical characteristic boundary.  In the explicit witness, the three
hyperplanes remain distinct.  Treating their intersection as a physical
causal structure therefore needs a separate source-operational quotient or
compatibility law.  Without that law, the envelope is mathematical control
data, not universal matter propagation.

This distinction is the exact reason to select a construction-or-obstruction
packet rather than claim success now.  The next packet must stress:
empty intersections, multiplicity and degeneracy, loss of strict margin,
independent sector variation, all-sector coverage, operational typing, and
hidden target fit.  It must return either one constructed proposal-only
candidate or one precise scoped obstruction.

\section{Source-extension classification}

The product and intersection operations are conservative mathematical
definitions once the sector polynomials are supplied.  The supply of
independently motivated sector principal data and any law giving the envelope
physical operational meaning are not derived by the current ontology.  They
are classified as a \emph{new ontology primitive candidate},
\emph{proposal-only}, and \emph{source-extension data}.  No adoption is
requested.

The missing source-side law for the selected packet is
\[
 \texttt{source\_operational\_common\_hyperbolicity\_envelope\_quotient\_law}.
\]
Current ontology does not derive that law.  Current adoption is blocked while
same-milestone research continuation remains open:
\texttt{blocked\_adoption\_open\_continuation}.

\section{Distance-to-GR matrix}

\begin{longtable}{>{\raggedright\arraybackslash}p{0.25\textwidth}>{\raggedright\arraybackslash}p{0.22\textwidth}>{\raggedright\arraybackslash}p{0.43\textwidth}}
\toprule
Burden & Status & Selector consequence\\
\midrule
Source ontology primitives & draft object exists & One proposal-only envelope packet is selected; no primitive is adopted.\\
Source equivalence EqSrc & Refuter stress passed & Unchanged.\\
RetainH & draft object exists & Unchanged.\\
GenH & draft object exists & Unchanged.\\
ObsLoc\(_{\rm lc}\) & blocked by missing primitive & Exact protocol typing remains restricted and constant without extension.\\
Resp\(_{\rm lc}\) & blocked by missing primitive & D7 is not evaluated.\\
\(M_{\rm src}\) & blocked by missing primitive & The smooth source patch remains proposal-only scaffolding.\\
\(g_{\rm eff}\) & blocked by missing primitive & The envelope packet is selected but not executed; no physical cone, scale, or metric exists.\\
Matter coupling & not started & A common time-covector domain would not establish universal matter propagation.\\
Einstein equations & not started & No action, variation, or field equation is supplied.\\
Finite-variation robustness & blocked by missing primitive & The next packet must test strict-margin robustness and empty intersections.\\
Benchmark promotion & human-gated & No benchmark action is attempted.\\
Gate Chair status & human-gated & No protected review is requested.\\
Current route freeze or hard-fail status & frozen negative & The shared-line strategy remains locally frozen; the selected envelope branch is materially distinct and unexecuted.\\
\bottomrule
\end{longtable}

\section{No-fog decision}

Exactly one result is produced: select a future Candidate Constructor packet
to build or precisely obstruct \(\Pprod\), \(\GammaInt\), and their proposed
source-operational quotient on an honestly restricted domain.  The selector
does not build that candidate.  It proves only that product hyperbolicity and
a nonempty common time-covector intersection are mathematically available in
a split-sector witness without equality of sector lines.

\begin{nonconclusion}
The common envelope is not yet a physical cone, conformal structure, metric,
clock, rod, universal propagation law, coupling, or Einstein equation.  The
sector principal data and operational quotient are not derived or adopted.
D7 adequacy is unevaluated, B2 is inactive, P4--T03 is locked, Distance-to-GR
is unchanged, and no Gate verdict, benchmark promotion, global no-go,
publication, push, external review, or completed derivation follows.
\end{nonconclusion}

\end{document}
